\section{评测体系设计}

\subsection{七臂总览}

评测套件按「每臂对准管线一段 + 一个组合层」设计：每条臂独立成 harness，产出机读落盘记录供计分卡与归因消费。评测器规格冻结件为 \texttt{docs/spec/benchmark.md}，共同底材为 corpus\_v3。表~\ref{tab:arms} 为总览。

\begin{table}[htbp]
\centering\small
\caption{七臂评测总览（B1--B7 + 辅助件）}
\label{tab:arms}
\setlength{\tabcolsep}{4pt}
\begin{tabular}{@{}lllll@{}}
\toprule
\# & 评测臂 & 测哪段 & 底材 & 核心指标 \\
\midrule
B1 & parsebench & 解析段 & corpus\_v3 + corpus39 & 解析成功/一致率/泄漏/死件/合平 \\
B2 & fixtures & 解析段（单元级） & 手构陷阱 \texttt{.tex} 集（10 件） & 陷阱断言通过率（T01--T98） \\
B3 & compilebench & 编译段 + 修复循环 & corpus\_v3 原始包 & 纯净/带瑕/失败、救回率、规则谱 \\
B4 & xlatbench & 翻译段 & corpus\_v3 译段抽样 & 硬契约率 / 延迟 / token 经济 \\
B4b & qualbench & 翻译质量 & e2e 产物段对 & LLM 评审 ESA 评分 / 缺陷分类 \\
B5 & e2ebench & 全链（组合） & corpus\_v3 子集 & 逐段漏斗 + 终态分布 + 引入回归 \\
B6 & validbench & 校验段 & 合成破坏案例 & 检出率 / 误报 / 延迟 \\
B7 & alignbench & 阅读体验 & 中英编译产物对 & 具名锚点保留率 / 链权 \\
\midrule
\multicolumn{5}{@{}l}{辅助件：stagerun（五段批量流水线）、gullet\_bench（展开机实测）、}\\
\multicolumn{5}{@{}l}{gate\_scorecard（M2 出口门）、translators\_bench（模拟臂工厂）、nightwatch}\\
\bottomrule
\end{tabular}
\end{table}

依赖序：corpus\_v3 $\to$ B1 $\to$ B4 $\to$ \{B3, B5\} $\to$ B7；B2/B6 以合成输入独立。

\subsection{评测现场：\texttt{bench/} 的构成}

所有评测资产集中在仓库 \texttt{bench/} 树，入口文档为 \texttt{bench/TIERS.md}（验证分层契约）与 \texttt{docs/spec/benchmark.md} §3（逐库横评协议——原 \texttt{bench/PROTOCOL.md} 已于 2026-09-20 退役）。后者定义「一个问题该在哪一级被抓住」：L0 单测/断言样例（秒级，CI 硬门）$\to$ L1 语料机制覆盖（分钟级）$\to$ L2 stagerun 子集回归（分钟到小时，机制台账定向重放）$\to$ L3 全层全臂过夜集成（里程碑门）；规则是「能低不高」——低层能钉住的问题不许只在批量上撞见，批量发现的真坑要沉淀回 L0 层断言。本报告 B1--B7 属 L1--L3 层读数。布局与职能：

\begin{table}[htbp]
\centering\small
\caption{\texttt{bench/} 现场布局（不熟悉仓库的读者对照此表读结果目录名）}
\label{tab:benchlayout}
\begin{tabular}{@{}l>{\raggedright\arraybackslash}p{10.8cm}@{}}
\toprule
路径 & 内容 \\
\midrule
\texttt{bench/py/} & 37 个评测脚本，分四簇——\textbf{批量驱动}：\texttt{stagerun.py} + \texttt{stage\_\{ingest,parse,xlat,compile,fixloop\}}（五段流水线，落盘记录逐格追加式写入）；\textbf{七臂评测驱动}：\texttt{parsebench}、\texttt{compilebench\_v3}、\texttt{xlatbench}、\texttt{qualbench}、\texttt{e2e\_mock\_bench}、\texttt{e2e\_real\_bench}、\texttt{fixloop\_bench}、\texttt{alignbench}、\texttt{validbench}、\texttt{gullet\_bench}、\texttt{fixture\_assert}；\textbf{门禁监控}：\texttt{gate\_scorecard}（M2 出口门）、\texttt{status\_panel}、\texttt{nightwatch}、\texttt{rundiff}、\texttt{triage}、\texttt{wave}、\texttt{stage\_timing}；\textbf{支撑件}：\texttt{benchlib} 公共库、\texttt{translators\_bench}（模拟译文工厂）、\texttt{corpus/} 语料构建管线、\texttt{iclr\_*} 对照语料渠、\texttt{report/} 一次性审计、\texttt{scratch/} 一次性探针 \\
\texttt{bench/ts/} & JS 侧三库横评驱动（latex-utensils/unified-latex/tree-sitter-latex），D0 选型证据产地 \\
\texttt{bench/corpus*} & 语料族：\texttt{corpus} 39 篇手挑陷阱、\texttt{corpus\_v2} 139 篇分层、\texttt{corpus\_v3} 13{,}266 篇八层钉版（\S\ref{sec:corpus}）、\texttt{corpus\_daily} RSS 日更渠（$\sim$1{,}200 篇/日）、\texttt{corpus\_iclr*} ICLR 对照渠 \\
\texttt{bench/fixtures/} & 10 件手构陷阱 \texttt{.tex}（\texttt{@Tnn}/\texttt{@Wnn}/\texttt{@Xn} 断言标记：T 为结构陷阱、W 为野外发现机制、X 为翻译层；逐字节即语义，不做任何格式化） \\
\texttt{bench/results/} & 评测产物区：每次运行一个日期目录（377+ 轮累积），含 records/summary/jsonl/PDF 中间件；大体量不入档，本报告读数全部由其抽取进 \texttt{data/} \\
\texttt{bench/work\_*/} & 编译/fixloop 逐格工作区（解压源树、aux、PDF；重产物 gitignored） \\
\bottomrule
\end{tabular}
\end{table}

\subsection{语料分层构造}
\label{sec:corpus}

corpus\_v3 是钉版批量语料：每格记录 \texttt{(channel, item, member, blob\_sha256)} 四元组钉住字节级来源（\texttt{resolved\_version=null}），落地为 \texttt{\{id\}/} 目录（\texttt{meta.json}、\texttt{raw.*}、\texttt{extracted/}）。四条渠道各有明确分工：\texttt{ia} = IA \texttt{arxiv-bulk} 月度块（a--d 时代带主力，覆盖至 2020-10）；\texttt{tiger} = HF \texttt{TIGER-Lab/arxiv-latex-5T}（e 带主力，成员截止 2412 月）；\texttt{arxiv\_eprint} = 单发源包取源（走产品 \texttt{acquire\_source}，补 TIGER 截止后 2501+ 盲区）；\texttt{hf\_scholarweave} = scholarweave 脱水通道（\texttt{figures\_stripped} 有损源——解析/翻译段可用，编译段因缺图源排除）。全库 13{,}266 格、原始包 18.8GB（含 \texttt{extracted/} 解压树全库约 46GB）。八层构成见图~\ref{fig:corpus}、分层定位见表~\ref{tab:layers}：核心/补强/扩展/热层为早期层，2026-09-19 评测/开发分轨扩层后收至 13{,}266 篇。

\begin{figure}[htbp]
\centering
\begin{tikzpicture}
\begin{axis}[
  xbar stacked,
  width=13cm, height=4.6cm,
  xmin=0, xmax=14000,
  scaled x ticks=false,
  xlabel={千篇}, ytick=data,
  xtick={0,2000,4000,6000,8000,10000,12000,14000},
  xticklabels={0,2,4,6,8,10,12,14},
  yticklabels={corpus\_v3 八层},
  legend style={at={(0.5,-0.28)},anchor=north,legend columns=4,font=\scriptsize},
  bar width=8mm,
  x label style={font=\small},
]
\addplot+[fill=cA!75,draw=cA] coordinates {(1000,0)};
\addplot+[fill=cB!75,draw=cB] coordinates {(200,0)};
\addplot+[fill=cD!75,draw=cD] coordinates {(3866,0)};
\addplot+[fill=cE!75,draw=cE] coordinates {(166,0)};
\addplot+[fill=cF!75,draw=cF] coordinates {(2000,0)};
\addplot+[fill=cC!75,draw=cC] coordinates {(1514,0)};
\addplot+[fill=cB!40,draw=cB] coordinates {(1500,0)};
\addplot+[fill=black!35,draw=black!60] coordinates {(3020,0)};
\legend{core 1000, booster 200, expand 3866, hot 166, dev\_vol 2000, dev\_recent 1514, dev\_failmine 1500, holdout 3020}
\end{axis}
\end{tikzpicture}
\caption{corpus\_v3 八层构成（13{,}266 篇，原始包 18.8GB / 含解压树约 46GB）。holdout 层为 EVAL\_ONLY 治理——\texttt{dev\_layers()} 实测排除，仅评测可读；dev\_* 三层与 expand 为开发期增量；core 为月度簇分层随机（30 簇，old 200/new 800）。}
\label{fig:corpus}
\end{figure}

\begin{table}[htbp]
\centering\small
\caption{corpus\_v3 八层定位与抽样框（渠道缩写：ia=IA 月块 / tiger=HF arxiv-latex-5T / eprint=单发取源 / sw=scholarweave 脱水）}
\label{tab:layers}
\begin{tabular}{@{}ll>{\raggedright\arraybackslash}p{3.4cm}>{\raggedright\arraybackslash}p{6.4cm}@{}}
\toprule
层 & 篇数 & 定位 & 渠道构成与抽样框 \\
\midrule
core & 1{,}000 & 核心随机层，评测主干 & ia 800 + tiger 200；30 个 IA 月簇分层随机（old 200 / new 800） \\
booster & 200 & 机制策展层 & ia 162 + tiger 38；野外发现机制提名制（107/109 W exemplars，\texttt{nominations/} 留审计轨迹） \\
expand & 3{,}866 & 体量扩展层 & ia 3{,}103 + tiger 697 + eprint 66；bulk 池 measure-then-sample \\
hot & 166 & 近期高引层 & arxiv\_eprint 166；OpenAlex 高引近期，专打 2501+ 盲区 \\
dev\_vol & 2{,}000 & 开发体量层 & ia 1{,}627 + tiger 373；3{,}671 tex $\cdot$ 2.76GB \\
dev\_recent & 1{,}514 & 近期开发层 & sw 1{,}065 + eprint 449；TIGER 截止后盲区双通道 \\
dev\_failmine & 1{,}500 & 机制挖掘层 & ia 1{,}500；FLAG\_RX 机制旗标定向挖 IA 94--15 年档陷阱 \\
holdout & 3{,}020 & 留出评测层 & ia 2{,}159 + tiger 540 + eprint 321；EVAL\_ONLY，同构造规则独立抽样 \\
\bottomrule
\end{tabular}
\end{table}

抽样框是「时代带 + 机制旗标带」的二维结构：五个时代带 a\_pre2007 / b\_2007\_11 / c\_2012\_16 / d\_2017\_20 / e\_2021\_25 各 $\sim$1{,}700--2{,}050 篇均匀打底（时间跨度 1994--2026），上面叠机制旗标带定向覆盖已知陷阱族（\texttt{flag:deadpkg} 600、\texttt{flag:docstyle209} 300、\texttt{flag:epsfig}、\texttt{flag:babel}、\texttt{flag:pdftex}、\texttt{flag:pstricks} 等若干）与 \texttt{sw|yymm} 月度带 1{,}065。学科面随 arXiv 主力盘分布（图~\ref{fig:corpusdist} 左）：hep-phys 2{,}961、math 2{,}418、cs 2{,}393、astro-ph 1{,}826、cond-mat 1{,}729 为五大盘。

物理形态与工程结构：tar 包 11{,}097、gz 单文件 2{,}144、single\_gz 12、stub 13（B07 有意构造的残缺件）；单篇文件数中位 7 / p90 28 / 最大 3{,}047（一篇论文三千文件），\texttt{.tex} 均 2.2 个，多文件工程 2{,}115 篇（16\%）——合平与多文档路由不是边角料而是主场景之一。许可面：arxiv-nonexclusive 6{,}786、missing 3{,}335、CC 系合计 1{,}187。全库 70 个聚簇键（\texttt{cluster\_id}，月簇/渠道混合）支撑簇稳健自举；4{,}472 格带 B01--B07/W/T 机制标签——即「评测发现 $\to$ 归因 $\to$ 策展回填」反馈环的载体。

\begin{figure}[htbp]
\centering
\begin{tikzpicture}
\begin{groupplot}[
  group style={group size=2 by 1,horizontal sep=18mm},
]
\nextgroupplot[
  xbar, bar width=3.4mm,
  width=7.6cm, height=5.2cm,
  xmin=0, xmax=3450,
  symbolic y coords={nucl,quant-ph,eess-stat,cond-mat,astro-ph,cs,math,hep-phys},
  ytick=data, y tick label style={font=\scriptsize},
  xlabel={篇数}, x label style={font=\small},
  nodes near coords, nodes near coords style={font=\tiny},
]
\addplot+[fill=cA!70,draw=cA] coordinates {(282,nucl) (552,quant-ph) (583,eess-stat) (1729,cond-mat) (1826,astro-ph) (2393,cs) (2418,math) (2961,hep-phys)};
\nextgroupplot[
  ybar, bar width=4.4mm,
  width=7.2cm, height=5.2cm,
  ymin=0, enlarge y limits={upper,value=0.22},
  symbolic x coords={a$<$07,b 07--11,c 12--16,d 17--20,e 21--25,sw,flag},
  xtick=data, x tick label style={font=\tiny,rotate=28},
  ylabel={篇数}, y label style={font=\small},
  nodes near coords, nodes near coords style={font=\tiny},
]
\addplot+[fill=cD!70,draw=cD] coordinates {(a$<$07,2039) (b 07--11,2047) (c 12--16,1936) (d 17--20,1726) (e 21--25,1817) (sw,1065) (flag,2636)};
\end{groupplot}
\end{tikzpicture}
\caption{corpus\_v3 构成两视图。左：学科组分布 top8（其余 470 篇未绘）；右：stratum 带分布——a--e 五时代带共 9{,}565 篇均匀打底，sw 月带 1{,}065，flag 机制带/dev 补/holdout recent 等合计 2{,}636。}
\label{fig:corpusdist}
\end{figure}

治理规则：holdout 层（3{,}020 篇）登记 \texttt{EVAL\_ONLY\_LAYERS}，所有开发期采样经 \texttt{dev\_layers()} 过滤；语料扩展只加层不删层；格数据不入库（gitignored），入库物为清单八件与选样报告（\texttt{selection\_report.md}）；机制台账 \texttt{mechanisms.jsonl}、B04/B06 宇宙×语料覆盖簿记 \texttt{eval\_coverage.json}、补强层提名审计轨迹 \texttt{nominations/} 均为 \texttt{bench/corpus/} 下正档——台账是「评测发现 $\to$ 归因 $\to$ 规则沉淀」反馈环的载体。

\subsection{逐臂详解}
\label{sec:arms}

\subsubsection{B1 parsebench —— 解析正确性}

\textbf{测什么}：\texttt{texlate.latex} 半解析器在真实语料上的正确性——能不能零崩溃跑完、能不能逐字节还原、可译译段里有没有漏进受保护内容。

\textbf{怎么测}：逐文件 30 秒超时，先定位主文件，再做多文件合平，随后调用解析器，最后跑重建三连测——逐字节一致（strict identity）、空译文假翻译（fake-translate，注入死占位符/孤儿探针）、占位译文（泄漏探针）。逐文件结果写入 \texttt{files.jsonl}，逐篇聚合成 \texttt{papers.json}，同时产出路由标签（reject / xelatex / minted / non-utf8 / no-hyperref——这组标签同时充当 B3 的静态路由金标准）以及合平覆盖率、多文档/无根标记。统计口径为核心层加权池化率与宏平均双报，置信区间用 Wilson 加月簇稳健自举；另有 oracle 对拍探针（plasTeX 规范化文本流对译段的字符级召回），仅作开发期召回探针、不进门槛。所有异常逐条人工归因后写回机制台账。

\textbf{当前读数}：28{,}904 文件解析成功率 100.0\%、逐字节一致率 99.99\%、泄漏率 0.004\%（图~\ref{fig:parse}、表~\ref{tab:snapshot}）。

\subsubsection{B2 fixtures —— 机制断言集}

\textbf{测什么}：parsebench 测分布，断言集测机制——「这个具体的坑处理了吗」逐条断言。

\textbf{底材}：10 件手构 \texttt{.tex}（tricky 主集 + 209 旧格式件 + multi 多文件工程件 + w/w73/wenc 野外机制件 + xlat-traps 翻译层件），\texttt{@Tnn}（结构陷阱）/\texttt{@Wnn}（野外发现机制）/\texttt{@Xn}（翻译层）三族标记；\textbf{逐字节即语义}，永不格式化。\textbf{生长机制}：parsebench 归因与台账中 \texttt{found-in-wild} 条目达到 covered 状态后提取固化为断言样例——语料里每个真坑沉淀为一条永久断言，断言面只增不减。

\textbf{当前读数}：98 条断言 pytest 全绿（L0 层 CI 硬门，断言函数与 \texttt{fixture\_assert.py} 共享）；BUG1 占位符尾回归断言实测为 0。

\subsubsection{B3 compilebench —— 编译与修复循环救回}

\textbf{测什么}：规范化 + ctex 注入 + 引擎 + 修复循环的真实救回能力——hjfy 约五千篇人肉沉淀护城河的对应物。

\textbf{怎么测}：论文 $\times$ 条件 $\times$ 引擎三维网格——条件取原文裸编（base）/ 中文注入（zh）/ 模拟译文（mock）三种，引擎取 xelatex 与 tectonic 两种；编译前先做路由预检（\texttt{documentstyle}$\to$直接拒编、eps/pstricks$\to$xelatex、minted-frozencache$\to$tectonic）。每格走「规范化 $\to$ 中文注入 $\to$ 沙箱编译（至多 2 轮、240 秒限时）$\to$ 修复循环 $\to$ 纯净判定三件套」；逐格判定（verdict）连同每轮 \texttt{\{cat,pay,rule,result\}} 动作落盘，直接灌回机制台账。失败归因按首错签名聚类（\texttt{verdict\_sig} 口径：当构成中计数最高的错误类超过首错类时改挂该众数类）。指标：纯净率（分条件/引擎/时代带）、失败$\to$出 PDF 救回率、规则触发/救回数、卡死率、修复轮数分布。

\textbf{门槛}：200 篇中文链条件编译成功率 $\geq$90\%（hjfy 的 95\% 为渐近线）、拒编判定正确率 100\%（与路由标签对拍）、无回归（原本纯净的格不被规则改脏）。

\textbf{当前读数}：500 格联合口径出 PDF 90.4\%（图~\ref{fig:compile}）；中文臂 xelatex 纯净率 +16.1 个百分点——规范化顺带修复源缺陷，产品链成本转正收益。

\subsubsection{B4 xlatbench / qualbench —— 翻译契约与质量}

\textbf{B4a 硬契约层}：译段 $\times$ 模型 $\times$ 重复次数网格，译段按类型分层抽样（正文段、图题、条目、脚注，外加压力样例——\verb|\bibitem| 前缀、\verb|\href| 等陷阱位）；每个响应过 L0 校验器，指标按序读——硬契约通过率（丢/造占位符、丢脆弱命令、校验错误）$\to$ 换序软信号（合法中文换序不降权）$\to$ 延迟 $\to$ 推理量 $\to$ token 经济。用途是模型选型与白名单刷新（免费集推广期到期即重跑）以及提示词措辞回归（每次提升 \texttt{prompt\_version} 必跑）。

\textbf{B4b qualbench 质量层}：LLM 评审（esa2 协议）逐段对打 0--100 分，并按错误分类 $\times$ 严重度三级打标（术语/流畅/准确/约定/漏译五族九类）；初裁 swe-2-max，分歧判例记为争议，由 swe-2-high 复核仲裁；产出均分/中位分、分数段分布、按译段类型与按论文的细分。门槛 M1：100 篇真实翻译端到端 $\geq$85\%。

\textbf{当前读数}：硬契约通过率 93.7\%（254/271）、ESA 均分 94.0 / 争议率 5.5\%（图~\ref{fig:qual}）。

\subsubsection{B5 e2ebench —— 全链组合}

\textbf{测什么}：单段绿不等于组合绿——逐环节成功率漏斗加终态分布。四模式：\textbf{Mode A} 位置忠实的模拟译文（用占位译文走全链验证机械链路：能否出 PDF、逐字节一致、零残留占位符；已知盲区是模拟翻译只认 ASCII 散文段，非 ASCII 散文会原样回显、过估忠实度——真实翻译臂不受影响）；\textbf{Mode B} 幻觉模拟（注入占位符丢失/幻觉破坏，验证校验链兜底——132 处破坏编译前 132/132 捕获，「出 PDF $\neq$ 成功」的实证来源）；\textbf{Mode C} 位置扰动模拟（随机移动约 10\% 占位符，量化回填鲁棒性）；\textbf{Mode D} 真实翻译接入（产品 SLA 观测点）。另有 \textbf{pipe-fix 救回臂}：把 pipe-xel 的产物树送进修复循环复判，联合口径取两臂较优者；\texttt{onfail} 语义经冒烟校准为「失败 + misschar/error 级带瑕」才救（带瑕已出 PDF，引擎重跑只会丢 PDF 救不了告警），注入拒编不救。

\textbf{当前读数}：模拟臂管线 xel 33/39 纯净 vs 裸编 19/39（管线净正 +14 格）；真实臂 60 篇联合口径出 PDF 96.7\%、回填残留 0、0 管线引入回归。

\subsubsection{B6 validbench —— 校验器自检}

\textbf{测什么}：L0/L1 校验器对 LLM 破坏的检出能力——校验器本身必须被评测（实现期真抓到过校验器自身 bug）。

\textbf{怎么测}：变异器对干净译段施加 10 类破坏（丢右括号、丢美元符、\verb|\end| 改名或删除、丢/拼错/凭空多占位符、\verb|\[| 不配对、幻觉宏、删引用键），在占位符层/原文层两层生成 \texttt{cases.jsonl}；另加对抗性手工探针（丢 \verb|\emph| 组、占位符合法换序、\verb|\cite {k}| 带空格、漏引用键、全角占位符等——验证合法改写不误报）。指标：检出率（按破坏类分列命中规则）、干净对误报率、仅告警率、每对延迟、编辑距离 $\leq$2 修复建议率；L1 tree-sitter 层同口径。B4 产出的真实失败案例反向喂回 B6，沉淀为新破坏类。

\textbf{门槛}：10 类破坏 100\% 检出、干净对 0 误报、L0 $\leq$1ms/对。

\textbf{当前读数}：corpus\_v2 1{,}503 对全绿（摊薄 0.374ms/对）——本批未重跑，列缺口清单第 7 项。

\subsubsection{B7 alignbench —— 锚点保留}

\textbf{测什么}：中文重编译后 PDF 具名锚点（named-destination）保留率——双语滚动同步的质量上限。

\textbf{怎么测}：用 pypdf 提取中英两侧具名锚点，同名锚点配对，计算保留率与最大权值单调链权重（节标题锚权 12、图锚 10、公式锚 4、引用锚 2，\texttt{page.*} 页锚权 0）。未用 hyperref 的工程（语料中约 31\%）走退化路径、只断言不崩溃；\textbf{保留率 $<$95\% 本身可当中文编译完整性探针}（ctex 共享计数器改 theorem 锚名的问题即由此捕获）。

\textbf{门槛}：有 hyperref 对保留率 $\geq$95\%；无 hyperref 对退化路径不崩。

\textbf{当前读数}：812 对门全过但覆盖薄——765 对无具名锚点，有效读数仅 42 个 hyperref 对（保留率中位 1.0），口径待复核（缺口清单第 3 项）。

\subsubsection{辅助件}

stagerun 五段批量流水线（逐格追加式落盘、按 \texttt{(id,arm,upstream)} 断点续跑、网关信号量限流），配 \texttt{gate\_scorecard} 做 M2 出口门（联合口径 = 编译 $\cup$ 修复两段逐格取优）；\texttt{translators\_bench} 为模拟臂译文工厂（忠实模拟 / 注入破坏 b、c / 位置扰动）；\texttt{gullet\_bench} 为展开机实测（200 篇全展开中位 39.6ms/19 步）；另有 \texttt{triage}（落盘记录 $\to$ 问题工单聚类）、\texttt{status\_panel}、\texttt{nightwatch}、\texttt{rundiff}、\texttt{wave}、\texttt{stage\_timing}、\texttt{preflight\_batch} 等监控与批前闸件。

\subsection{门槛体系}

各臂门槛（\texttt{docs/spec/benchmark.md} §8 门槛汇总）：parsebench——解析成功率 100\%、逐字节一致率 $\geq 99.5\%$、泄漏率 $\leq 0.15\%$、死占位符与孤儿译段 $=0$、合平覆盖率 $\geq 99\%$；M2 出口门（stagerun 计分卡）——联合口径出 PDF 率 $\geq 90\%$、纯净率 $\geq 90\%$。统计口径：核心层加权池化率 + 宏平均双报，置信区间用 Wilson 加月簇稳健自举。
